%% =====================================================================
%%  Capítulo 4 -- Representación en espacio de estado
%%  Migrado al formato nuevo (noviembre de 2026) desde Handout3_Ctrl.tex.
%%  Correcciones de fondo: ver MIGRACION.md y la caja "Correcciones" de la
%%  sección 4.8 (levitador sin la constante C, historia del RIAS, ejemplo del
%%  determinante, funciones de transferencia del bioreactor, truco de canon).
%% =====================================================================

\capitulo[cap:ee]{Representación en espacio de estado}
  {Otra consecuencia es que se está impulsando que la educación sea vista como una capacitación que solo conduce al trabajo, y no como un medio que conduce al pensamiento crítico y a la transformación.}
  {Junot Díaz}

\begin{objetivos}[partes/04-espacio-estado.html]
  \item Definir el estado de un sistema y escribir su representación en espacio de estado, en particular en variables de fase.
  \item Pasar de una función de transferencia a una representación de estado, y de vuelta, y relacionar los valores propios de $\Am$ con los polos.
  \item Resolver la ecuación de estado con la matriz de transición $\Phi(t)=e^{\Am t}$.
  \item Linealizar un modelo no lineal de varias variables alrededor de un equilibrio y escribirlo en espacio de estado.
\end{objetivos}

%% ---------------------------------------------------------------------
\section{Dos maneras de describir un sistema}
\label{sec:ee-intro}
%% ---------------------------------------------------------------------

Hasta aquí hemos descrito los sistemas con \textbf{funciones de transferencia}: una relación algebraica entre la entrada y la salida, en el dominio de Laplace. Es el método \emph{clásico} o \emph{frecuencial}, y tiene ventajas grandes: interconectar subsistemas es multiplicar y sumar bloques (\cref{cap:so}), y la estabilidad y el transitorio se leen casi de inmediato en los polos. Su límite es igual de claro: solo sirve para sistemas lineales e invariantes en el tiempo, o para lo que se pueda aproximar así, y solo describe la relación entre una entrada y una salida, no lo que pasa por dentro.

La \textbf{representación en espacio de estado} describe el sistema con un conjunto de ecuaciones diferenciales de \emph{primer orden}, una por cada variable que hace falta para conocer su evolución. Es más general (admite sistemas no lineales, variantes en el tiempo y con varias entradas y salidas) y muestra la dinámica interna. En este curso la usamos para sistemas lineales, en tres momentos: aquí, para escribir y resolver modelos; en el \cref{cap:es}, para la estabilidad; y en el \cref{cap:re}, para diseñar controladores que usan todo el estado.

\begin{historia}[Del Sputnik al Apolo]
El lanzamiento del Sputnik, en 1957, despertó en Occidente el interés por lo que hacía la ingeniería soviética. El primer congreso mundial de la IFAC se reunió en Moscú en 1960, y allí muchos ingenieros de control constataron que la escuela soviética trabajaba de otra manera: describía los procesos con ecuaciones diferenciales en lugar de funciones de transferencia, analizaba la estabilidad con los métodos de Lyapunov en lugar de los de Bode y Nyquist, y estudiaba el desempeño óptimo con el principio del máximo de Pontryagin en lugar del método de Wiener-Hopf. Así empezó lo que se llamó \textbf{control moderno}.

En Estados Unidos las ideas se difundieron desde la Universidad de Columbia, en el grupo de John R. Ragazzini, donde se formó Rudolf E. Kalman, y desde el RIAS (\emph{Research Institute for Advanced Studies}) de Baltimore, cuyo centro de matemáticas dirigía Solomon Lefschetz y donde Kalman trabajó desde 1958. En el propio congreso de Moscú, Kalman presentó una teoría general de los sistemas de control en espacio de estado~\citep{Kal60}; con J.~E. Bertram aplicó el método de Lyapunov, y con Richard Bucy extendió la teoría de filtrado de Wiener a sistemas variantes en el tiempo~\citep{KaBu61}. En 1963 mostró con precisión cómo se relacionan la función de transferencia y la representación de estado: qué partes del sistema ve la primera y cuáles le quedan ocultas~\citep{Kal63}. El filtro de Kalman voló en la navegación del programa Apolo.
\end{historia}

%% ---------------------------------------------------------------------
\section{Un ejemplo: el ciclo global del carbono}
\label{sec:ee-carbono}
%% ---------------------------------------------------------------------

Antes de definir nada, un modelo. Li, Jarvis y Leedal describen el ciclo global del carbono con cinco reservorios que intercambian carbono entre sí~\citep{LJL09} (\cref{fig:ee-carbono}). Las variables, en gigatoneladas de carbono (GtC), son las masas en la atmósfera $m_A$, la vegetación $m_V$, el suelo $m_S$, la superficie del océano $m_{SO}$ y el océano profundo $m_{DO}$; la entrada es la emisión humana $u_E$, en GtC/año. Cada flujo es proporcional a la masa del reservorio de donde sale, con un coeficiente de intercambio $k$ en 1/año.

\begin{figure}[ht]
\centering
\includegraphics[width=0.85\linewidth]{cap04/carbono}
\caption{Ciclo global del carbono con cinco reservorios~\citep{LJL09}. Cada flecha es un flujo proporcional a la masa del reservorio de origen.}
\label{fig:ee-carbono}
\end{figure}

Un balance de masa por reservorio (lo que entra menos lo que sale) da
\begin{align*}
  \dot m_A    &= u_E-(k_{O1}+k_{L1})m_A+k_{L2}m_V+k_{L4}m_S+k_{O2}m_{SO},\\
  \dot m_V    &= k_{L1}m_A-(k_{L2}+k_{L3})m_V,\\
  \dot m_S    &= k_{L3}m_V-k_{L4}m_S,\\
  \dot m_{SO} &= k_{O1}m_A-(k_{O2}+k_{O3})m_{SO}+k_{O4}m_{DO},\\
  \dot m_{DO} &= k_{O3}m_{SO}-k_{O4}m_{DO}.
\end{align*}
Son cinco ecuaciones de primer orden, y se pueden escribir de una vez en forma matricial, $\dot x=\Am x+\Bm u_E$, con $x=[m_A\;\,m_V\;\,m_S\;\,m_{SO}\;\,m_{DO}]^\top$,
\begin{equation}
  \Am=\begin{bmatrix}
  -(k_{O1}+k_{L1}) & k_{L2} & k_{L4} & k_{O2} & 0\\
  k_{L1} & -(k_{L2}+k_{L3}) & 0 & 0 & 0\\
  0 & k_{L3} & -k_{L4} & 0 & 0\\
  k_{O1} & 0 & 0 & -(k_{O2}+k_{O3}) & k_{O4}\\
  0 & 0 & 0 & k_{O3} & -k_{O4}
  \end{bmatrix},\qquad
  \Bm=\begin{bmatrix}1\\0\\0\\0\\0\end{bmatrix}.
  \label{eq:ee-carbono}
\end{equation}

\begin{ejemplo}[label=ej:ee-carbono]{Lo que dice la estructura de $\Am$}
Con $k_{L1}=0.0842$, $k_{L2}=0.0684$, $k_{L3}=0.0159$, $k_{L4}=0.0332$, $k_{O1}=0.0846$, $k_{O2}=0.0433$, $k_{O3}=0.0011$ y $k_{O4}=0.0001$ (todos en 1/año), ¿qué pasa con $u_E=8$ GtC/año constante?
\solucion Cada columna de $\Am$ suma cero: lo que sale de un reservorio entra a otro. Por eso $\mathbf 1^\top\Am=0$ y, como $\mathbf 1^\top\Bm=1$,
\[
  \deriv{}{t}\bigl(m_A+m_V+m_S+m_{SO}+m_{DO}\bigr)=\mathbf 1^\top\dot x=u_E :
\]
el carbono total crece exactamente al ritmo de las emisiones; los intercambios solo lo redistribuyen. La misma propiedad obliga a que $\Am$ tenga un valor propio en $\lambda=0$ (sus filas son linealmente dependientes, porque suman cero). Los otros cuatro son $-0.227$, $-0.067$, $-0.036$ y $-0.00058$ por año: constantes de tiempo de $4.4$, $15$, $28$ y $1729$ años. Con $u_E=8$ GtC/año durante 500 años, la atmósfera pasa de $590$ a $1833$ GtC (\cref{fig:ee-carbono-sim}): el valor propio nulo actúa como un integrador, y con emisión constante no hay estado estacionario.
\end{ejemplo}

\begin{figure}[ht]
\centering
\includegraphics[width=\linewidth]{cap04/f_carbono}
\caption{Ejemplo~\ref{ej:ee-carbono}: 500 años de emisiones constantes de 8 GtC/año, desde $x(0)=[590\;\,550\;\,1500\;\,900\;\,37000]^\top$ GtC.}
\label{fig:ee-carbono-sim}
\end{figure}

\codigorepo{cap04_espacio_estado/C4_Ej5_Carbono.m}{matrices}{las matrices, la conservación de la masa y los valores propios}

Con este ejemplo en mente veamos qué es el estado, cómo se escoge y qué se gana con él.

%% ---------------------------------------------------------------------
\section{El estado de un sistema}
\label{sec:ee-estado}
%% ---------------------------------------------------------------------

La premisa de la mecánica de Newton es que la evolución futura de un sistema queda determinada por su condición presente y por las fuerzas que actúen sobre él. Esa condición presente es el estado.

\begin{definicion}[label=def:ee-estado]{Estado}
El \textbf{estado} de un sistema dinámico es el conjunto más pequeño de variables cuyo valor en un instante $t_0$, junto con las entradas para $t\ge t_0$, determina por completo el comportamiento del sistema para todo $t\ge t_0$~\citep{Fri86}. Esas variables son las \textbf{variables de estado}; agrupadas forman el \textbf{vector de estado} $x(t)$, y el espacio $n$-dimensional cuyos ejes son las variables de estado es el \textbf{espacio de estado}.
\end{definicion}

El número de variables de estado (el \textbf{orden} del sistema) es el número de condiciones iniciales que hacen falta para resolver sus ecuaciones; es único. Las variables, en cambio, no lo son: el mismo sistema admite muchas elecciones distintas. Una regla práctica: en un sistema físico, cada elemento que almacena energía de manera independiente aporta una variable de estado (la posición de un resorte, la velocidad de una masa, la carga de un condensador, la corriente de una bobina, el nivel de un tanque).

\paragraph{La forma general.} El modelo se escribe como $n$ ecuaciones de primer orden, las \textbf{ecuaciones de estado}, y una \textbf{ecuación de salida} algebraica:
\begin{equation}
  \dot x=f(x,u,t),\qquad y=h(x,u,t),
  \label{eq:ee-no-lineal}
\end{equation}
con $x\in\mathbb R^n$, $u\in\mathbb R^m$ y $y\in\mathbb R^p$. Si el sistema es lineal e invariante en el tiempo,
\begin{equation}
  \boxed{\;\dot x=\Am x+\Bm u,\qquad y=\Cm x+\Dm u\;}
  \label{eq:ee-lineal}
\end{equation}
donde $\Am$ ($n\times n$) es la matriz del sistema, $\Bm$ ($n\times m$) la de entrada, $\Cm$ ($p\times n$) la de salida y $\Dm$ ($p\times m$) la de transmisión directa. Este capítulo trabaja con \eqref{eq:ee-lineal}, salvo en la \cref{sec:ee-linealizacion}, donde se llega a ella desde \eqref{eq:ee-no-lineal}.

\subsection{Variables de fase}

Si la ecuación diferencial no tiene derivadas de la entrada,
\begin{equation}
  y^{(n)}+a_{n-1}y^{(n-1)}+\dots+a_1\dot y+a_0y=b_0u,
  \label{eq:ee-edo}
\end{equation}
la elección más directa es tomar la salida y sus $n-1$ derivadas: las \textbf{variables de fase},
\[
  x_1=y,\quad x_2=\dot y,\quad\dots,\quad x_n=y^{(n-1)} .
\]
Cada variable es la derivada de la anterior, y la última ecuación sale de \eqref{eq:ee-edo}:
\begin{equation}
  \dot x=\underbrace{\begin{bmatrix}0&1&0&\cdots&0\\0&0&1&\cdots&0\\\vdots&&&\ddots&\vdots\\0&0&0&\cdots&1\\-a_0&-a_1&-a_2&\cdots&-a_{n-1}\end{bmatrix}}_{\Am}x
  +\underbrace{\begin{bmatrix}0\\0\\\vdots\\0\\b_0\end{bmatrix}}_{\Bm}u,
  \qquad y=\underbrace{\begin{bmatrix}1&0&\cdots&0\end{bmatrix}}_{\Cm}x .
  \label{eq:ee-fase}
\end{equation}
Los coeficientes del polinomio característico aparecen, cambiados de signo, en la última fila de $\Am$.

\begin{ejemplo}[label=ej:ee-masa]{Masa, resorte y amortiguador}
Para el sistema de la \cref{fig:ee-masa}, con fricción viscosa $b\dot y$ y resorte lineal $ky$, la segunda ley de Newton da $M\ddot y+b\dot y+ky=u$. Escriba su modelo de estado.
\solucion Es de segundo orden: dos variables de estado. Con $x_1=y$ (posición) y $x_2=\dot y$ (velocidad),
\[
  \dot x_1=x_2,\qquad \dot x_2=-\frac kMx_1-\frac bMx_2+\frac1Mu,\qquad y=x_1,
\]
es decir,
\[
  \Am=\begin{bmatrix}0&1\\-\dfrac kM&-\dfrac bM\end{bmatrix},\quad
  \Bm=\begin{bmatrix}0\\\dfrac1M\end{bmatrix},\quad
  \Cm=\begin{bmatrix}1&0\end{bmatrix},\quad \Dm=0 .
\]
Con $M=2$, $k=16$ y un pulso de $5$\,N durante $20$\,s, la \cref{fig:ee-masa-sim} muestra el efecto de $b$: con $b=2$ los valores propios de $\Am$ son $-0.5\pm2.78j$ y la masa oscila; con $b=16$ son reales ($-1.17$ y $-6.83$) y no hay oscilación. En los cuatro casos, la posición se asienta en $5/16=0.3125$\,m mientras dura el pulso.
\end{ejemplo}

\begin{figure}[ht]
\centering
\includegraphics[width=0.55\linewidth]{cap04/masa_resorte_u}
\caption{Masa, resorte y amortiguador: entrada $u$ (fuerza), salida $y$ (posición).}
\label{fig:ee-masa}
\end{figure}

\begin{figure}[ht]
\centering
\includegraphics[width=0.95\linewidth]{cap04/f_masa_resorte}
\caption{Ejemplo~\ref{ej:ee-masa}: posición y velocidad para cuatro valores de $b$, con $M=2$ y $k=16$.}
\label{fig:ee-masa-sim}
\end{figure}

\codigorepo{cap04_espacio_estado/C4_Ej1_MasaResorte.m}{fase}{el modelo de estado de la masa-resorte}
\codigorepo{cap04_espacio_estado/C4_Ej1_MasaResorte.m}{simulacion}{la simulación con \cmd{lsim} para cuatro valores de $b$}

\paragraph{El computador analógico.} Una ventaja práctica del modelo de estado es que se arma con tres tipos de bloques: \textbf{integradores}, \textbf{sumadores} y \textbf{ganancias}. Así se simulaban los sistemas en los computadores analógicos, antes de los digitales: cada variable de estado es la salida de un integrador, y su condición inicial se fijaba con un potenciómetro. La \cref{fig:ee-analogico} muestra la masa-resorte armada de esa manera. Simulink conserva la misma lógica.

\begin{figure}[ht]
\centering
\includegraphics[width=0.85\linewidth]{cap04/computador_analogico}
\caption{La masa-resorte como en un computador analógico: dos integradores en cascada, cuyas salidas son las variables de estado.}
\label{fig:ee-analogico}
\end{figure}

\begin{modelo}[Modelo de Simulink: \texttt{crearModelo\_MasaResorte}]
\texttt{codigo/cap04\_espacio\_estado/crearModelo\_MasaResorte.m} arma la \cref{fig:ee-analogico} con bloques \emph{Integrator}, \emph{Sum} y \emph{Gain}, y un pulso de $5$\,N durante $20$\,s. Cambie \texttt{b} en el espacio de trabajo y vuelva a simular para reproducir la \cref{fig:ee-masa-sim}.
\end{modelo}

\begin{ejemplo}[label=ej:ee-dos-masas]{Dos masas}
En la \cref{fig:ee-dos-masas}, $M_1$ está unida a la pared por un amortiguador $b$ y a $M_2$ por un resorte $k$; la fuerza $u$ actúa sobre $M_2$. Las ecuaciones de Newton son
\[
  M_1\ddot y_1+b\dot y_1+k(y_1-y_2)=0,\qquad M_2\ddot y_2+k(y_2-y_1)=u .
\]
¿Cuántas variables de estado hacen falta? Escriba el modelo si se miden las dos posiciones.
\solucion Las posiciones $y_1$ y $y_2$ no bastan: hay dos ecuaciones de segundo orden y hacen falta cuatro condiciones iniciales. Con $x_1=y_1$, $x_2=y_2$, $x_3=\dot y_1$ y $x_4=\dot y_2$,
\[
  \Am=\begin{bmatrix}0&0&1&0\\0&0&0&1\\-\dfrac k{M_1}&\dfrac k{M_1}&-\dfrac b{M_1}&0\\\dfrac k{M_2}&-\dfrac k{M_2}&0&0\end{bmatrix},\quad
  \Bm=\begin{bmatrix}0\\0\\0\\\dfrac1{M_2}\end{bmatrix},\quad
  \Cm=\begin{bmatrix}1&0&0&0\\0&1&0&0\end{bmatrix},\quad\Dm=0 .
\]
Aquí $p=2$: dos salidas.
\end{ejemplo}

\begin{figure}[ht]
\centering
\includegraphics[width=0.7\linewidth]{cap04/dos_masas}
\caption{Ejemplo~\ref{ej:ee-dos-masas}: dos masas unidas por un resorte.}
\label{fig:ee-dos-masas}
\end{figure}

\begin{pausa}
Un circuito RLC en serie, alimentado por una fuente de voltaje, ¿cuántas variables de estado tiene? Proponga dos elecciones distintas.
\respuesta{Dos, porque hay dos elementos que almacenan energía de forma independiente. Por ejemplo, el voltaje del condensador y la corriente de la bobina; o la carga del condensador y su derivada (que es la corriente), que son variables de fase.}
\end{pausa}

%% ---------------------------------------------------------------------
\section{De la función de transferencia al estado}
\label{sec:ee-ft-estado}
%% ---------------------------------------------------------------------

Una función de transferencia sin ceros,
\[
  \frac{Y(s)}{U(s)}=\frac{b_0}{s^n+a_{n-1}s^{n-1}+\dots+a_1s+a_0},
\]
es la transformada de \eqref{eq:ee-edo} con condiciones iniciales nulas, y las variables de fase \eqref{eq:ee-fase} salen de inmediato.

\begin{ejemplo}[label=ej:ee-fase]{Una función de transferencia de tercer orden}
Escriba en variables de fase $\dfrac{Y(s)}{U(s)}=\dfrac{24}{s^3+9s^2+26s+24}$.
\solucion La ecuación diferencial es $\dddot y+9\ddot y+26\dot y+24y=24u$. Con $x_1=y$, $x_2=\dot y$ y $x_3=\ddot y$,
\[
  \Am=\begin{bmatrix}0&1&0\\0&0&1\\-24&-26&-9\end{bmatrix},\quad
  \Bm=\begin{bmatrix}0\\0\\24\end{bmatrix},\quad
  \Cm=\begin{bmatrix}1&0&0\end{bmatrix},\quad\Dm=0 .
\]
Los valores propios de $\Am$ son $-2$, $-3$ y $-4$: los polos de la función de transferencia.
\end{ejemplo}

\subsection{Cuando hay ceros}

Si el numerador tiene grado $m\ge1$, la ecuación diferencial tiene derivadas de la entrada y la salida y sus derivadas ya no sirven como estado. El truco es partir la función de transferencia en dos (\cref{fig:ee-intermedia}): primero el denominador, que produce una variable intermedia $Z(s)$ sin derivadas de la entrada,
\[
  Z(s)=\frac{1}{s^n+a_{n-1}s^{n-1}+\dots+a_0}\,U(s),
\]
a la que sí se le aplican las variables de fase ($x_1=z$, $x_2=\dot z$, \dots); y después el numerador, que solo combina esas variables:
\[
  Y(s)=(b_ms^m+\dots+b_1s+b_0)Z(s)\quad\Longrightarrow\quad y=b_0x_1+b_1x_2+\dots+b_mx_{m+1}.
\]
Resultado: la misma $\Am$ de \eqref{eq:ee-fase}, $\Bm=[0\;\cdots\;0\;\,1]^\top$ y $\Cm=[b_0\;\,b_1\;\cdots\;b_m\;\,0\;\cdots\;0]$. Los ceros quedan todos en $\Cm$. Esta es la forma que entrega la función del curso \cmd{variablesFase}, y la que se usa en el \cref{cap:re}. (Para $m=0$ coincide con \eqref{eq:ee-fase} salvo en dónde se pone $b_0$: en $\Bm$ o en $\Cm$.)

\begin{figure}[ht]
\centering
\includegraphics[width=0.9\linewidth]{cap04/variable_intermedia}
\caption{(a) Función de transferencia con ceros. (b) Partición con la variable intermedia $Z(s)$: las variables de fase se definen sobre $z$, y el numerador solo forma la salida.}
\label{fig:ee-intermedia}
\end{figure}

\begin{ejemplo}[label=ej:ee-ceros]{Con ceros}
Escriba en variables de fase $\dfrac{Y(s)}{U(s)}=\dfrac{s^2+7s+2}{s^3+9s^2+26s+24}$.
\solucion La misma $\Am$ del ejemplo~\ref{ej:ee-fase}, $\Bm=[0\;\,0\;\,1]^\top$ y $\Cm=[2\;\,7\;\,1]$: $y=2z+7\dot z+\ddot z$.
\end{ejemplo}

\codigorepo{cap04_espacio_estado/C4_Ej2_FTaEstado.m}{fase}{ejemplo~\ref{ej:ee-fase}}
\codigorepo{cap04_espacio_estado/C4_Ej2_FTaEstado.m}{numerador}{ejemplo~\ref{ej:ee-ceros} con la función \cmd{variablesFase}}

\paragraph{Lo que hace \matlab.} El comando \cmd{tf2ss} no devuelve las variables de fase. Para el ejemplo~\ref{ej:ee-ceros} entrega
\[
  \Am_1=\begin{bmatrix}-9&-26&-24\\1&0&0\\0&1&0\end{bmatrix},\quad
  \Bm_1=\begin{bmatrix}1\\0\\0\end{bmatrix},\quad
  \Cm_1=\begin{bmatrix}1&7&2\end{bmatrix} :
\]
las mismas variables, pero numeradas al revés ($x_1$ de \matlab{} es nuestro $x_3$). Invertir el orden de los estados con $\mathbb P=$\cmd{fliplr(eye(3))} recupera nuestra forma: $\mathbb P\Am_1\mathbb P$, $\mathbb P\Bm_1$, $\Cm_1\mathbb P$. Las dos representaciones dan la misma función de transferencia: \alerta{la representación de estado no es única, la función de transferencia sí}. La \cref{sec:ee-estado-ft} explica por qué.

\codigorepo{cap04_espacio_estado/C4_Ej2_FTaEstado.m}{matlab}{\cmd{tf2ss} y el cambio de orden de los estados}

\begin{pausa}
¿Qué pasa con $\Dm$ si el numerador tiene el mismo grado que el denominador, por ejemplo $\dfrac{2s+3}{s+1}$?
\respuesta{Primero se divide: $\dfrac{2s+3}{s+1}=2+\dfrac{1}{s+1}$. La parte constante va directo de la entrada a la salida: $\Dm=2$; el resto, $\Am=-1$, $\Bm=1$, $\Cm=1$.}
\end{pausa}

%% ---------------------------------------------------------------------
\section{Solución de la ecuación de estado}
\label{sec:ee-solucion}
%% ---------------------------------------------------------------------

\paragraph{Respuesta libre.} Empecemos sin entrada: $\dot x=\Am x$, con $x(0)$ conocido. Integrando a ambos lados, $x(t)=x(0)+\int_0^t\Am x(\tau)\,\dif\tau$; reemplazando $x(\tau)$ por esa misma expresión una y otra vez,
\[
  x(t)=\Bigl(\Id+\Am t+\frac{\Am^2t^2}{2!}+\frac{\Am^3t^3}{3!}+\cdots\Bigr)x(0).
\]
La serie es la de la exponencial, ahora con una matriz:
\begin{equation}
  x(t)=e^{\Am t}x(0)=\Phi(t)\,x(0) .
  \label{eq:ee-libre}
\end{equation}

\begin{definicion}[label=def:ee-transicion]{Matriz de transición de estado}
$\Phi(t)=e^{\Am t}=\displaystyle\sum_{k=0}^{\infty}\frac{\Am^kt^k}{k!}$ es la \textbf{matriz de transición de estado}: lleva el estado del instante $0$ al instante $t$. Cumple
\[
  \Phi(0)=\Id,\qquad \Phi^{-1}(t)=\Phi(-t),\qquad \Phi(t_1+t_2)=\Phi(t_1)\Phi(t_2),\qquad \Phi^k(t)=\Phi(kt).
\]
\end{definicion}

Calcular $e^{\Am t}$ no es trivial: Moler y Van Loan reunieron diecinueve maneras de hacerlo, todas con algún defecto numérico~\citep{MoLo78}. A mano, la más útil es la de Laplace (\cref{sec:ee-estado-ft}); en \matlab, \cmd{expm}.

\paragraph{Respuesta forzada.} Con entrada, $\dot x-\Am x=\Bm u$. Multiplicando por $e^{-\Am t}$, el lado izquierdo es $\deriv{}{t}\bigl(e^{-\Am t}x\bigr)$, y al integrar de $0$ a $t$:
\begin{equation}
  x(t)=\underbrace{e^{\Am t}x(0)}_{\text{respuesta libre}}+\underbrace{\int_0^te^{\Am(t-\tau)}\Bm\,u(\tau)\,\dif\tau}_{\text{respuesta forzada}},
  \qquad y(t)=\Cm x(t)+\Dm u(t).
  \label{eq:ee-forzada}
\end{equation}
Es la versión matricial de lo visto en el \cref{cap:po}: la respuesta total es la suma de la respuesta a las condiciones iniciales y la convolución de la entrada con la respuesta al impulso.

\begin{ejemplo}[label=ej:ee-transicion]{Matriz de transición de un segundo orden}
Con $\Am=\begin{bmatrix}0&1\\-2&-3\end{bmatrix}$ y $\Bm=[0\;\,1]^\top$: halle $\Phi(t)$, la respuesta al escalón unitario con $x(0)=0$ y la respuesta libre con $x(0)=[1\;\,{-1}]^\top$.
\solucion Con el procedimiento de la \cref{sec:ee-estado-ft},
\[
  \Phi(t)=\begin{bmatrix}2e^{-t}-e^{-2t} & e^{-t}-e^{-2t}\\-2e^{-t}+2e^{-2t} & -e^{-t}+2e^{-2t}\end{bmatrix},
  \qquad \Phi(0.5)=\begin{bmatrix}0.8452&0.2387\\-0.4773&0.1292\end{bmatrix}.
\]
Con \eqref{eq:ee-forzada} y $u=1$: $x_1(t)=\int_0^t\bigl(e^{-(t-\tau)}-e^{-2(t-\tau)}\bigr)\dif\tau=\frac12-e^{-t}+\frac12e^{-2t}$, y $x_2=\dot x_1=e^{-t}-e^{-2t}$. Con $x(0)=[1\;\,{-1}]^\top$ y sin entrada, $x(t)=e^{-t}[1\;\,{-1}]^\top$: un solo modo, porque esa condición inicial es un vector propio de $\Am$ (con valor propio $-1$).
\end{ejemplo}

\codigorepo{cap04_espacio_estado/C4_Ej3_Transicion.m}{simbolico}{$\Phi(t)$ con la resolvente y \cmd{ilaplace}}
\codigorepo{cap04_espacio_estado/C4_Ej3_Transicion.m}{numerico}{\cmd{expm} y las propiedades de $\Phi$}
\codigorepo{cap04_espacio_estado/C4_Ej3_Transicion.m}{forzada}{la respuesta forzada contra \cmd{lsim}}

\begin{pausa}
Si $x(0)=[1\;\,{-2}]^\top$ en el ejemplo~\ref{ej:ee-transicion}, ¿qué modo aparece?
\respuesta{Solo $e^{-2t}$: $[1\;\,{-2}]^\top$ es vector propio de $\Am$ con valor propio $-2$ ($\Am[1\;\,{-2}]^\top=[-2\;\,4]^\top$). Entonces $x(t)=e^{-2t}[1\;\,{-2}]^\top$.}
\end{pausa}

%% ---------------------------------------------------------------------
\section{Del estado a la función de transferencia}
\label{sec:ee-estado-ft}
%% ---------------------------------------------------------------------

Transformando $\dot x=\Am x+\Bm u$ con condiciones iniciales nulas, $sX(s)=\Am X(s)+\Bm U(s)$, de donde $X(s)=(s\Id-\Am)^{-1}\Bm U(s)$, y con la ecuación de salida,
\begin{equation}
  \boxed{\;G(s)=\frac{Y(s)}{U(s)}=\Cm(s\Id-\Am)^{-1}\Bm+\Dm\;}
  \label{eq:ee-ss2tf}
\end{equation}
siempre que $s\Id-\Am$ sea invertible. Como $(s\Id-\Am)^{-1}=\operatorname{adj}(s\Id-\Am)/\det(s\Id-\Am)$, cada entrada de $G(s)$ es un cociente de polinomios con denominador $\det(s\Id-\Am)$, el \textbf{polinomio característico} de $\Am$. Por eso los polos de $G(s)$ están entre los valores propios de $\Am$.

La matriz $(s\Id-\Am)^{-1}$ se llama \textbf{resolvente}, y es la transformada de la matriz de transición: transformando la serie de $e^{\Am t}$ término a término, $\Lap{e^{\Am t}}=\frac{\Id}{s}+\frac{\Am}{s^2}+\frac{\Am^2}{s^3}+\cdots$, y al multiplicar por $(s\Id-\Am)$ todo se cancela salvo $\Id$. Entonces
\begin{equation}
  \Phi(t)=e^{\Am t}=\iLap{(s\Id-\Am)^{-1}} .
  \label{eq:ee-resolvente}
\end{equation}

\begin{receta}[Matriz de transición por Laplace]
\begin{enumerate}
  \item Forme $s\Id-\Am$ y calcule su determinante, el polinomio característico.
  \item Inviértala: $(s\Id-\Am)^{-1}=\dfrac{\operatorname{adj}(s\Id-\Am)}{\det(s\Id-\Am)}$. Para $2\times2$, $\begin{bmatrix}a&b\\c&d\end{bmatrix}^{-1}=\dfrac{1}{ad-bc}\begin{bmatrix}d&-b\\-c&a\end{bmatrix}$.
  \item Antitransforme cada elemento, con fracciones parciales.
\end{enumerate}
\end{receta}

\paragraph{Adjunta y determinante.} La adjunta es la transpuesta de la matriz de cofactores: $\operatorname{adj}(\Am)=[C_{ij}]^\top$, con $C_{ij}=(-1)^{i+j}M_{ij}$ y $M_{ij}$ el determinante de la matriz que queda al quitar la fila $i$ y la columna $j$ (el \emph{menor}). El determinante se puede desarrollar por cualquier fila $i$ o columna $k$: $\det\Am=\sum_ja_{ij}C_{ij}=\sum_ia_{ik}C_{ik}$. Conviene escoger la fila o la columna con más ceros.

\begin{ejemplo}[label=ej:ee-det]{Desarrollo por cofactores}
Calcule el determinante de $\begin{bmatrix}2&4&1\\3&0&2\\2&0&3\end{bmatrix}$.
\solucion La segunda columna tiene dos ceros. Los menores son $M_{12}=\begin{vmatrix}3&2\\2&3\end{vmatrix}=5$, $M_{22}=\begin{vmatrix}2&1\\2&3\end{vmatrix}=4$ y $M_{32}=\begin{vmatrix}2&1\\3&2\end{vmatrix}=1$, y los cofactores, $C_{12}=-5$, $C_{22}=4$ y $C_{32}=-1$. Entonces $\det=4\cdot C_{12}+0\cdot C_{22}+0\cdot C_{32}=-20$.
\end{ejemplo}

\begin{ejemplo}[label=ej:ee-motor]{Motor de corriente continua}
Un motor con inercia se modela como $\dot x_1=x_2$, $\dot x_2=-\alpha x_2+\beta u$, $y=x_1$, con $\alpha,\beta>0$~\citep{Fri86}. Halle $\Phi(t)$ y $G(s)$.
\solucion $s\Id-\Am=\begin{bmatrix}s&-1\\0&s+\alpha\end{bmatrix}$, $\det=s(s+\alpha)$ y
\[
  (s\Id-\Am)^{-1}=\begin{bmatrix}\dfrac1s&\dfrac{1}{s(s+\alpha)}\\[2mm]0&\dfrac{1}{s+\alpha}\end{bmatrix},
  \qquad
  \Phi(t)=\begin{bmatrix}1&\dfrac{1-e^{-\alpha t}}{\alpha}\\0&e^{-\alpha t}\end{bmatrix}.
\]
Con \eqref{eq:ee-ss2tf}, $G(s)=\dfrac{\beta}{s(s+\alpha)}$: el valor propio en cero es el integrador que convierte velocidad en posición.
\end{ejemplo}

\subsection{¿Valores propios o polos?}

Los polos de $G(s)$ son valores propios de $\Am$, pero \alerta{no siempre todos los valores propios son polos}: puede haber cancelaciones entre el numerador y $\det(s\Id-\Am)$.

\begin{ejemplo}[label=ej:ee-cancelacion]{Una función de transferencia que esconde modos}
Con $\Am=\begin{bmatrix}0&1&0\\0&0&1\\-6&-11&-6\end{bmatrix}$ y $\Bm=[0\;\,0\;\,1]^\top$, compare $\Cm_a=[1\;\,0\;\,0]$ y $\Cm_b=[2\;\,3\;\,1]$.
\solucion Los valores propios son $-1$, $-2$ y $-3$. Con $\Cm_a$, $G_a=\dfrac{1}{(s+1)(s+2)(s+3)}$. Con $\Cm_b$, el numerador es $s^2+3s+2=(s+1)(s+2)$ y
\[
  G_b(s)=\frac{(s+1)(s+2)}{(s+1)(s+2)(s+3)}=\frac{1}{s+3}:
\]
un solo polo para tres valores propios. Los modos $e^{-t}$ y $e^{-2t}$ siguen en el estado (aparecen si $x(0)\neq0$), pero no llegan a la salida. Aquí son estables y la cancelación es inofensiva; si uno de ellos fuera inestable, la función de transferencia lo escondería. Es la pregunta que responde la \textbf{observabilidad}, en el \cref{cap:re}.
\end{ejemplo}

\codigorepo{cap04_espacio_estado/C4_Ej4_EstadoaFT.m}{ss2tf}{\cmd{ss2tf}, \cmd{eig} y \cmd{dcgain}}
\codigorepo{cap04_espacio_estado/C4_Ej4_EstadoaFT.m}{cancelacion}{ejemplo~\ref{ej:ee-cancelacion}}

\paragraph{Cambio de variables.} Si $x=\mathbb T\bar x$ con $\mathbb T$ invertible, el sistema en las nuevas variables es $\bar\Am=\mathbb T^{-1}\Am\mathbb T$, $\bar\Bm=\mathbb T^{-1}\Bm$, $\bar\Cm=\Cm\mathbb T$, $\bar\Dm=\Dm$, y
\[
  \bar\Cm(s\Id-\bar\Am)^{-1}\bar\Bm=\Cm\mathbb T\bigl(\mathbb T^{-1}(s\Id-\Am)\mathbb T\bigr)^{-1}\mathbb T^{-1}\Bm=\Cm(s\Id-\Am)^{-1}\Bm .
\]
Por eso hay infinitas representaciones de estado con la misma función de transferencia, y los valores propios no cambian. Si las columnas de $\mathbb T$ son vectores propios de $\Am$ (con valores propios distintos), $\bar\Am$ es diagonal, cada estado nuevo es un \textbf{modo} desacoplado, y $\Phi(t)=\mathbb T\,\operatorname{diag}(e^{\lambda_it})\,\mathbb T^{-1}$.

\codigorepo{cap04_espacio_estado/C4_Ej4_EstadoaFT.m}{similitud}{la forma diagonal y $\Phi(t)$}

\begin{ideaclave}
\textbf{Lo que es único y lo que no.}\enspace La función de transferencia y los valores propios de $\Am$ no dependen de cómo se escojan las variables de estado; las matrices sí. Y la función de transferencia puede tener menos polos que valores propios: describe solo la parte del sistema que conecta la entrada con la salida.
\end{ideaclave}

%% ---------------------------------------------------------------------
\section{Linealización en espacio de estado}
\label{sec:ee-linealizacion}
%% ---------------------------------------------------------------------

En la \cref{sec:so-linealizacion} linealizamos ecuaciones de una variable. La representación de estado permite hacerlo con todas a la vez. Sea $\dot x=f(x,u)$, $y=h(x,u)$, y sea $(\eq x,\eq u)$ un equilibrio: $f(\eq x,\eq u)=0$. Con las variables desviación $\dev x=x-\eq x$, $\dev u=u-\eq u$ y $\dev y=y-h(\eq x,\eq u)$, la serie de Taylor de primer orden de cada componente da
\begin{equation}
  \dot{\dev x}=\Am\dev x+\Bm\dev u,\qquad \dev y=\Cm\dev x+\Dm\dev u,
  \label{eq:ee-lineal-desv}
\end{equation}
donde cada matriz es un \textbf{jacobiano} evaluado en el equilibrio:
\begin{equation}
  a_{ij}=\left.\frac{\partial f_i}{\partial x_j}\right|_{\eq x,\eq u},\quad
  b_{ij}=\left.\frac{\partial f_i}{\partial u_j}\right|_{\eq x,\eq u},\quad
  c_{ij}=\left.\frac{\partial h_i}{\partial x_j}\right|_{\eq x,\eq u},\quad
  d_{ij}=\left.\frac{\partial h_i}{\partial u_j}\right|_{\eq x,\eq u}.
  \label{eq:ee-jacobianos}
\end{equation}
Cada equilibrio tiene su propio modelo lineal, y en general dinámicas distintas.

\begin{ejemplo}[label=ej:ee-duffing]{La ecuación de Duffing}
Para $\ddot x+\delta\dot x+x^3-x=0$ con $\delta>0$: escriba el modelo de estado, halle los equilibrios y linealice alrededor de cada uno.
\solucion Con $x_1=x$, $x_2=\dot x$: $\dot x_1=x_2$, $\dot x_2=x_1-x_1^3-\delta x_2$. Los equilibrios cumplen $x_2=0$ y $x_1(1-x_1^2)=0$: son $(0,0)$ y $(\pm1,0)$. El jacobiano es
\[
  J(x_1,x_2)=\begin{bmatrix}0&1\\1-3x_1^2&-\delta\end{bmatrix} .
\]
En el origen, $\lambda^2+\delta\lambda-1=0$ y $\lambda=\frac{-\delta\pm\sqrt{\delta^2+4}}{2}$: un valor propio positivo y uno negativo para \emph{todo} $\delta>0$ (su producto es $-1$). Es un \textbf{punto de silla}, inestable. En $(\pm1,0)$, $\lambda^2+\delta\lambda+2=0$: $\omega_n=\sqrt2$ y $\zeta=\delta/(2\sqrt2)$, estables, subamortiguados si $\delta<\sqrt8=2.83$. Con $\delta=0.3$: $0.861$ y $-1.161$ en el origen, y $-0.15\pm1.41j$ en $(\pm1,0)$. La \cref{fig:ee-duffing} muestra el retrato de fase: cada condición inicial termina en uno de los dos pozos, y la frontera entre ellos (en negro) pasa por el origen en la dirección del vector propio estable de la silla. Ningún modelo lineal único reproduce eso.
\end{ejemplo}

\begin{figure}[ht]
\centering
\includegraphics[width=0.6\linewidth]{cap04/f_duffing}
\caption{Ejemplo~\ref{ej:ee-duffing} con $\delta=0.3$: retrato de fase, la frontera entre los dos dominios de atracción (negro) y tres trayectorias.}
\label{fig:ee-duffing}
\end{figure}

\begin{ejemplo}[label=ej:ee-levitador]{Levitador magnético}
Un electroimán sostiene una esfera de acero de masa $M$ (\cref{fig:ee-levitador}). La fuerza magnética es $Ci^2/y$, con $y$ la distancia de la esfera al imán~\citep{barie1996maglev}, y la bobina tiene resistencia $R$ e inductancia $L$:
\[
  M\ddot y=Mg-C\frac{i^2}{y},\qquad e=Ri+L\frac{\dif i}{\dif t}.
\]
Linealice alrededor de la posición $y=y_d$ con $M=0.05$\,kg, $R=10\,\Omega$, $L=0.5$\,H, $C=10^{-3}$\,N\,m/A$^2$, $g=9.81$\,m/s$^2$ y $y_d=0.01$\,m.
\solucion Con $x_1=y$, $x_2=\dot y$, $x_3=i$ y $u=e$:
\[
  \dot x_1=x_2,\qquad \dot x_2=g-\frac CM\frac{x_3^2}{x_1},\qquad \dot x_3=-\frac RLx_3+\frac1Lu .
\]
En el equilibrio, $\eq x_1=y_d$, $\eq x_2=0$, $\eq x_3=\sqrt{Mgy_d/C}=2.215$\,A y $\eq u=R\eq x_3=22.15$\,V. Los jacobianos son
\[
  \Am=\begin{bmatrix}0&1&0\\\dfrac CM\dfrac{(\eq x_3)^2}{(\eq x_1)^2}&0&-\dfrac{2C\eq x_3}{M\eq x_1}\\0&0&-\dfrac RL\end{bmatrix}
  =\begin{bmatrix}0&1&0\\981&0&-8.86\\0&0&-20\end{bmatrix},\qquad
  \Bm=\begin{bmatrix}0\\0\\2\end{bmatrix}.
\]
Note que $a_{21}=g/y_d$: en el equilibrio la fuerza magnética iguala al peso, y si la esfera se acerca, la fuerza crece más rápido que el peso. Los valores propios son $\pm31.3$ y $-20$: \alerta{el equilibrio es inestable}, y la función de transferencia $\dfrac{\dev Y(s)}{\dev E(s)}=\dfrac{-17.72}{(s-31.3)(s+31.3)(s+20)}$ lo confirma. La \cref{fig:ee-levitador-sim} compara el modelo lineal con el no lineal tras un desplazamiento inicial de $0.1$\,mm: coinciden durante unos $0.1$\,s (error del $4$\,\%) y luego se separan, cuando la esfera ya se alejó varios milímetros.
\end{ejemplo}

\begin{figure}[ht]
\centering
\includegraphics[width=0.6\linewidth]{cap04/levitador}
\caption{Levitador magnético: circuito de la bobina y fuerzas sobre la esfera.}
\label{fig:ee-levitador}
\end{figure}

\begin{figure}[ht]
\centering
\includegraphics[width=0.75\linewidth]{cap04/f_levitador}
\caption{Ejemplo~\ref{ej:ee-levitador}: la esfera parte $0.1$\,mm por debajo del equilibrio, con $u=\eq u$.}
\label{fig:ee-levitador-sim}
\end{figure}

\codigorepo{cap04_espacio_estado/C4_Ej6_Linealizacion.m}{duffing}{ejemplo~\ref{ej:ee-duffing}}
\codigorepo{cap04_espacio_estado/C4_Ej6_Linealizacion.m}{levitador}{ejemplo~\ref{ej:ee-levitador}}

\begin{ejemplo}[label=ej:ee-bioreactor]{Un bioreactor}
En un bioreactor, la biomasa $x_1$ (células) consume el sustrato $x_2$ (azúcar, desechos) y se multiplica. Un modelo simplificado es~\citep{Beq03}
\[
  \dot x_1=(\mu-u_1)x_1,\qquad \dot x_2=u_1(u_2-x_2)-\frac{\mu x_1}{Y},\qquad
  \mu(x_2)=\frac{\mu_{\max}x_2}{k_m+x_2+k_1x_2^2},
\]
donde $u_1$ es la tasa de dilución (caudal/volumen, la variable manipulada), $u_2$ la concentración de sustrato en la alimentación (una perturbación), $Y$ el rendimiento y $\mu$ la tasa de crecimiento con \emph{inhibición por sustrato} (con $k_1=0$ sería el modelo de Monod). Con $\mu_{\max}=0.53$\,h$^{-1}$, $k_m=0.12$\,g/L, $k_1=0.4545$\,L/g, $Y=0.4$, $\eq u_1=0.3$\,h$^{-1}$ y $\eq u_2=4$\,g/L, halle los equilibrios y linealice.
\solucion De $\dot x_1=0$: o $x_1=0$, o $\mu(x_2)=u_1$. El primer caso da el \emph{lavado}, $(0,\,4)$: las células se fueron del reactor. El segundo da $0.1362x_2^2-0.23x_2+0.036=0$, con $x_2=0.1746$ y $x_2=1.5122$, y de $\dot x_2=0$, $x_1=Y(u_2-x_2)$: los equilibrios $(1.5302,\,0.1746)$ y $(0.9951,\,1.5122)$. Con $\mu=u_1$ en el equilibrio, los jacobianos son
\[
  \Am=\begin{bmatrix}0&\eq x_1\mu'\\-\dfrac{u_1}{Y}&-u_1-\dfrac{\eq x_1\mu'}{Y}\end{bmatrix},\qquad
  \Bm=\begin{bmatrix}-\eq x_1&0\\u_2-\eq x_2&u_1\end{bmatrix},
\]
con $\mu'=\dif\mu/\dif x_2$. En $(0.9951,\,1.5122)$, $\mu'=-0.068$ y
\[
  \Am=\begin{bmatrix}0&-0.0679\\-0.75&-0.1302\end{bmatrix},\qquad
  \Bm=\begin{bmatrix}-0.9951&0\\2.4878&0.3\end{bmatrix},
\]
con valores propios $0.170$ y $-0.3$: \alerta{inestable}. En $(1.5302,\,0.1746)$, $\mu'=+0.591$ y los valores propios son $-0.3$ y $-2.26$: estable. La diferencia está en el signo de $\mu'$, que se anula en $x_2=\sqrt{k_m/k_1}=0.514$: a la derecha de ese máximo, más sustrato \emph{frena} el crecimiento (\cref{fig:ee-bioreactor}). Con las dos salidas $y=x$, las funciones de transferencia en el punto de operación son
\begin{align*}
  \frac{\dev X_1}{\dev U_1}&=\frac{-0.9951s-0.2985}{s^2+0.1302s-0.0510}, &
  \frac{\dev X_1}{\dev U_2}&=\frac{-0.0204}{s^2+0.1302s-0.0510},\\
  \frac{\dev X_2}{\dev U_1}&=\frac{2.488s+0.7463}{s^2+0.1302s-0.0510}, &
  \frac{\dev X_2}{\dev U_2}&=\frac{0.3\,s}{s^2+0.1302s-0.0510}.
\end{align*}
\end{ejemplo}

\begin{figure}[ht]
\centering
\includegraphics[width=\linewidth]{cap04/f_bioreactor}
\caption{Ejemplo~\ref{ej:ee-bioreactor}. Izquierda: la tasa de crecimiento corta a $u_1=0.3$ en dos puntos. Derecha: desde $(1,1)$ y desde cerca del equilibrio inestable, las trayectorias van al equilibrio estable o al lavado $(0,4)$.}
\label{fig:ee-bioreactor}
\end{figure}

\codigorepo{cap04_espacio_estado/C4_Ej7_Bioreactor.m}{equilibrios}{los tres equilibrios}
\codigorepo{cap04_espacio_estado/C4_Ej7_Bioreactor.m}{jacobiano}{jacobianos, valores propios y funciones de transferencia}

\begin{modelo}[Modelo de Simulink: \texttt{crearModelo\_Bioreactor}]
\texttt{codigo/cap04\_espacio\_estado/crearModelo\_Bioreactor.m} arma el modelo no lineal con dos entradas y los dos estados como salidas. Con él, \cmd{trim} busca el equilibrio para las entradas fijas y \cmd{linmod} entrega $\Am$, $\Bm$, $\Cm$ y $\Dm$ del modelo linealizado (las instrucciones están en el script \texttt{C4\_Ej7\_Bioreactor}).
\end{modelo}

\begin{cuidado}[El modelo lineal vale cerca del equilibrio]
Un modelo linealizado describe las desviaciones \emph{pequeñas}. Qué tan pequeñas, depende del sistema: en el levitador, unas décimas de milímetro; en el bioreactor, la linealización alrededor de un equilibrio no dice nada sobre el lavado. Y un equilibrio inestable en lazo abierto, como el del levitador o el del bioreactor, es justamente el que pide control.
\end{cuidado}

%% ---------------------------------------------------------------------
\section{Para llevar}
\label{sec:ee-cierre}
%% ---------------------------------------------------------------------

\begin{parallevar}
\begin{enumerate}
  \item El estado es el conjunto mínimo de variables que, con las entradas futuras, determina todo el comportamiento futuro. Su número (el orden) es único; las variables no.
  \item Modelo lineal: $\dot x=\Am x+\Bm u$, $y=\Cm x+\Dm u$. En variables de fase, la última fila de $\Am$ lleva los coeficientes del polinomio característico con signo cambiado; si hay ceros, van en $\Cm$ (variable intermedia $Z$). \cmd{tf2ss} numera los estados al revés.
  \item Solución: respuesta libre más respuesta forzada, $x(t)=e^{\Am t}x(0)+\int_0^te^{\Am(t-\tau)}\Bm u(\tau)\,\dif\tau$. La matriz de transición es $\Phi(t)=e^{\Am t}=\iLap{(s\Id-\Am)^{-1}}$.
  \item Función de transferencia: $G(s)=\Cm(s\Id-\Am)^{-1}\Bm+\Dm$. Sus polos son valores propios de $\Am$, pero una cancelación puede esconder valores propios.
  \item Un cambio de variables $x=\mathbb T\bar x$ cambia las matrices pero no la función de transferencia ni los valores propios.
  \item Linealizar es evaluar jacobianos en un equilibrio. Cada equilibrio tiene su modelo, y el modelo solo vale cerca de él.
\end{enumerate}
\end{parallevar}

\begin{profundiza}[Correcciones respecto a la versión anterior del capítulo]
(i) En el levitador, la ecuación de estado había perdido la constante $C$: es $\dot x_2=g-\frac CM\frac{x_3^2}{x_1}$, el equilibrio es $\eq x_3=\sqrt{Mgy_d/C}$ y $a_{21}=\frac CM\frac{(\eq x_3)^2}{(\eq x_1)^2}=g/y_d$ (en la matriz faltaba además el cuadrado de $\eq x_3$). (ii) En el ejemplo del determinante, el segundo cofactor es $C_{22}=4$ (estaba rotulado $C_{12}$). (iii) En el bioreactor, la cuarta función de transferencia repetía la tercera (\texttt{numU2(1,:)} en lugar de \texttt{numU2(2,:)}); los números del \texttt{trim} correspondían a $k_1=0.454$ y no a $0.4545$, y se recalcularon todos con el jacobiano analítico. (iv) El truco con \cmd{canon(\ldots,'companion')} y trasposiciones para recuperar las variables de fase dependía de la versión de \matlab; se reemplazó por la inversión del orden de los estados. (v) Historia: Lefschetz no creó el RIAS, sino que dirigió su centro de matemáticas; Pontryagin aportó el principio del máximo (no el ``cálculo de variaciones''); la conexión entre función de transferencia y estado que mostró Kalman (1963) es la de qué partes del sistema ve cada descripción, no que el espacio de estado cubra todo sistema no lineal. (vi) Erratas: $k_{04}\to k_{O4}$, $x_1=y\to x_1=y_1$ en el ejemplo de las dos masas, y una nota tomada de Wikipedia se reemplazó por la definición. (vii) Se agregaron los valores numéricos de todos los ejemplos, las figuras y el retrato de fase de Duffing; el resumen del capítulo anterior que cerraba el handout se retiró (está en el \cref{cap:so}). (viii) Notación: $\bar x\to\eq x$ para el equilibrio y $\dev x$ para la desviación, como en el \cref{cap:so}.
\end{profundiza}

%% ---------------------------------------------------------------------
\section{Problemas}
\label{sec:ee-problemas}
%% ---------------------------------------------------------------------

Los problemas marcados con \etiquetaM{} requieren \matlab/\simulink; los marcados con \etiquetaEstrella{} son de profundización.

\begin{problema}[variables de fase]\label{prob:ee-fase}
Escriba en variables de fase, con $\Am$, $\Bm$, $\Cm$ y $\Dm$: (a) $\dfrac{10}{s^3+6s^2+11s+6}$; (b) $2\ddot y+6\dot y+8y=4u$; (c) $\dfrac{2s^2+7s+3}{s^3+5s^2+8s+4}$, con la variable intermedia $Z(s)$. (d) \etiquetaM{} Compare con \cmd{tf2ss} y recupere su forma invirtiendo el orden de los estados. (e) \etiquetaM{} Aplique \cmd{ss2tf} a las dos representaciones de (c). ¿Qué propiedad ilustra el resultado?
\end{problema}

\begin{problema}[resolvente y polos]\label{prob:ee-resolvente}
Para (i) $\Am=\begin{bmatrix}0&1\\-6&-5\end{bmatrix}$ y (ii) $\Am=\begin{bmatrix}-2&1&0\\0&-3&1\\0&0&-4\end{bmatrix}$, con $\Bm=[0\;\cdots\;0\;\,1]^\top$, $\Cm=[1\;\,0\;\cdots]$ y $\Dm=0$: (a) calcule $(s\Id-\Am)^{-1}$ con la adjunta; (b) halle $G(s)$, sus polos y ceros; (c) compare con los valores propios de $\Am$; (d) halle la ganancia DC de (i) y el valor final ante escalón unitario; (e) \etiquetaM{} verifique con \cmd{ss2tf}, \cmd{eig} y \cmd{dcgain}.
\end{problema}

\begin{problema}[matriz de transición]\label{prob:ee-transicion}
(a) Calcule $\Phi(t)$ para $\Am=\begin{bmatrix}0&1\\-2&-3\end{bmatrix}$ con los tres pasos de la receta. (b) \etiquetaM{} Verifique numéricamente, con $t_1=0.5$\,s y $t_2=1.2$\,s, que $\Phi(0)=\Id$, $\Phi^{-1}(t_1)=\Phi(-t_1)$ y $\Phi(t_1+t_2)=\Phi(t_1)\Phi(t_2)$. (c) Deduzca la $\Phi(t)$ del motor del ejemplo~\ref{ej:ee-motor}. (d) Con $\Bm=[0\;\,1]^\top$ y $x(0)=0$, halle $x(t)$ ante escalón unitario con \eqref{eq:ee-forzada} y verifique en Laplace. (e) \etiquetaM{} Con $x(0)=[1\;\,{-1}]^\top$ y entrada nula, compare con \cmd{expm} en varios instantes. ¿Por qué aparece un solo modo?
\end{problema}

\begin{problema}[levitador]\label{prob:ee-levitador}
Con el modelo y los datos del ejemplo~\ref{ej:ee-levitador}: (a) escriba el modelo de estado no lineal conservando $C$; (b) halle $(\eq x,\eq u)$; (c) obtenga $\Am$ y $\Bm$ y muestre que $a_{21}=g/y_d$; explique por qué; (d) halle los valores propios; ¿qué concluye? (e) obtenga $\dev Y(s)/\dev E(s)$ y compare sus polos con (d); (f) \etiquetaM{} simule el modelo no lineal desde el equilibrio desplazado $0.1$\,mm y compare con el lineal; ¿hasta cuándo es aceptable la predicción?
\end{problema}

\begin{problema}[ciclo del carbono]\label{prob:ee-carbono}
Con el modelo y los coeficientes del ejemplo~\ref{ej:ee-carbono}: (a) construya $\Am$ y $\Bm$; (b) demuestre que $\mathbf 1^\top\Am=0$ y que $\frac{\dif}{\dif t}(\mathbf 1^\top x)=u_E$; interprete; (c) deduzca, sin calcular, que $\lambda=0$ es valor propio, y \etiquetaM{} halle los cinco y las constantes de tiempo; (d) \etiquetaM{} simule 500 años con $u_E=8$ GtC/año desde $x(0)=[590\;\,550\;\,1500\;\,900\;\,37000]^\top$; ¿hay estado estacionario? (e) halle $M_A(s)/U_E(s)$ con $\Cm=[1\;\,0\;\,0\;\,0\;\,0]$ e identifique el polo responsable del crecimiento.
\end{problema}

\begin{problema}[bioreactor]\label{prob:ee-bioreactor}
Con el modelo del ejemplo~\ref{ej:ee-bioreactor}: (a) derive $\mu'(x_2)$, muestre que se anula en $x_2=\sqrt{k_m/k_1}$ y ubique ese valor respecto a los equilibrios; (b) calcule a mano las ocho derivadas de $\Am$ y $\Bm$ en $(0.9951,\,1.5122)$; (c) halle los valores propios; ¿cuál modo es inestable? (d) \etiquetaM{} obtenga las cuatro funciones de transferencia; (e) repita (b) y (c) en $(1.5302,\,0.1746)$ y explique la diferencia con el signo de $\mu'$; (f) \etiquetaM{} simule desde $x(0)=(1,1)$; ¿a qué equilibrio llega?
\end{problema}

\begin{problema}[valores propios y polos]\label{prob:ee-cancelacion}
Con $\Am$, $\Bm$, $\Cm_a$ y $\Cm_b$ del ejemplo~\ref{ej:ee-cancelacion}: (a) halle los valores propios y las dos funciones de transferencia; (b) explique el resultado con $\Cm_b$ y reformule con precisión la frase ``los valores propios de $\Am$ son los polos''; (c) \etiquetaM{} simule los dos sistemas con $x(0)=[1\;\,0\;\,0]^\top$ y entrada nula, y muestre que hay modos en el estado que no llegan a la salida; (d) con $\mathbb T$ formada por los vectores propios, obtenga la forma diagonal y verifique que $G$ no cambia; (e) demuestre que todo cambio de variables preserva $G$; (f) \etiquetaM{} verifique $\Phi(t)=\mathbb T\operatorname{diag}(e^{\lambda_it})\mathbb T^{-1}$ contra \cmd{expm}.
\end{problema}

\begin{problema}[Duffing]\label{prob:ee-duffing}
Con $\ddot x+\delta\dot x+x^3-x=0$, $\delta>0$: (a) muestre que el origen es una silla para todo $\delta>0$; (b) halle el $\delta$ a partir del cual $(\pm1,0)$ deja de ser subamortiguado, y clasifique esos equilibrios para $\delta=0.3$ y $\delta=3$; (c) para $\delta=0.3$, halle $\Phi(t)$ del modelo lineal en $(1,0)$ y su tiempo de establecimiento; (d) \etiquetaM{} simule desde $(0.9,0)$, $(0.5,0)$ y $(-0.1,\,0.05)$ y compare con el modelo lineal del equilibrio correspondiente; (e) \etiquetaEstrella{} construya el retrato de fase en $[-2,2]^2$ y trace la frontera entre los dominios de atracción; relaciónela con los vectores propios del origen; (f) ¿por qué ningún modelo lineal reproduce lo de (e)?
\end{problema}

%% ---------------------------------------------------------------------
\section{Soluciones seleccionadas}
%% ---------------------------------------------------------------------

\paragraph{Problema~\ref{prob:ee-fase}}
(a) $\Am=\begin{bmatrix}0&1&0\\0&0&1\\-6&-11&-6\end{bmatrix}$, $\Bm=[0\;\,0\;\,10]^\top$, $\Cm=[1\;\,0\;\,0]$. (b) Dividiendo por $2$: $\ddot y+3\dot y+4y=2u$, $\Am=\begin{bmatrix}0&1\\-4&-3\end{bmatrix}$, $\Bm=[0\;\,2]^\top$. (c) $\Am=\begin{bmatrix}0&1&0\\0&0&1\\-4&-8&-5\end{bmatrix}$, $\Bm=[0\;\,0\;\,1]^\top$, $\Cm=[3\;\,7\;\,2]$. (e) La misma $G(s)$: la representación no es única, la función de transferencia sí.

\paragraph{Problema~\ref{prob:ee-resolvente}}
(i) $(s\Id-\Am)^{-1}=\dfrac{1}{s^2+5s+6}\begin{bmatrix}s+5&1\\-6&s\end{bmatrix}$, $G=\dfrac{1}{(s+2)(s+3)}$, sin ceros, ganancia DC $1/6$ y $y_\infty=1/6$. (ii) $G=\dfrac{1}{(s+2)(s+3)(s+4)}$ (los valores propios de una triangular son su diagonal), ganancia DC $1/24$.

\paragraph{Problema~\ref{prob:ee-transicion}}
(a) La del ejemplo~\ref{ej:ee-transicion}. (c) La del ejemplo~\ref{ej:ee-motor}. (d) $x_1=\frac12-e^{-t}+\frac12e^{-2t}$, $x_2=e^{-t}-e^{-2t}$; en Laplace, $X_1=\frac{1}{s(s+1)(s+2)}$. (e) $x(t)=e^{-t}[1\;\,{-1}]^\top$: la condición inicial es un vector propio.

\paragraph{Problema~\ref{prob:ee-levitador}}
(b) $\eq x=[0.01\;\,0\;\,2.215]^\top$, $\eq u=22.15$\,V. (c) $a_{21}=981$\,s$^{-2}$, $a_{23}=-8.86$, $a_{33}=-20$, $\Bm=[0\;\,0\;\,2]^\top$. (d) $\pm31.3$ y $-20$: inestable en lazo abierto. (e) $\dfrac{-17.72}{s^3+20s^2-981s-19620}$. (f) A $0.1$\,s el error del modelo lineal es de $4$\,\%; a $0.15$\,s, de $16$\,\%.

\paragraph{Problema~\ref{prob:ee-carbono}}
(c) $\lambda=0$, $-0.0006$, $-0.0355$, $-0.0673$ y $-0.2274$ por año; constantes de tiempo de $1729$, $28$, $15$ y $4.4$ años. (d) El total crece $8$ GtC/año; no hay estado estacionario. La atmósfera llega a $1833$ GtC en 500 años. (e) El polo en $s=0$.

\paragraph{Problema~\ref{prob:ee-bioreactor}}
(a) $\mu'=\mu_{\max}\dfrac{k_m-k_1x_2^2}{(k_m+x_2+k_1x_2^2)^2}$, nulo en $x_2=0.514$, entre los dos equilibrios. (c) $0.170$ y $-0.3$. (e) $\Am=\begin{bmatrix}0&0.905\\-0.75&-2.562\end{bmatrix}$, valores propios $-0.3$ y $-2.26$. (f) Llega a $(1.5302,\,0.1746)$.

\paragraph{Problema~\ref{prob:ee-cancelacion}}
(a) $-1$, $-2$, $-3$; $G_a=\frac{1}{(s+1)(s+2)(s+3)}$ y $G_b=\frac{1}{s+3}$. (b) Los polos de $G$ son valores propios de $\Am$; los valores propios que se cancelan no aparecen como polos.

\paragraph{Problema~\ref{prob:ee-duffing}}
(a) El producto de los valores propios es $\det J(0,0)=-1<0$. (b) $\delta=\sqrt8$; con $\delta=0.3$: $\zeta=0.106$, $\omega_n=1.41$, focos estables; con $\delta=3$: $\zeta=1.06$, valores propios $-1$ y $-2$, nodos estables. (c) $\ts\approx4/(\zeta\omega_n)=26.7$\,s. (d) $(0.9,0)$ y $(0.5,0)$ van a $(1,0)$; $(-0.1,\,0.05)$ va a $(-1,0)$: su componente sobre el vector propio inestable del origen es negativa.

\respuestaspausas

